\documentclass[10pt]{amsart}
\usepackage[a4paper,margin=23mm]{geometry}
\usepackage{amsmath,amssymb,amsthm,mathtools,hyperref,graphicx,xspace,letltxmacro,xpatch}
\usepackage[english]{babel}
\newtheorem{theo}{Theorem}[section]
\newtheorem{lemm}[theo]{Lemma}\newtheorem{prop}[theo]{Proposition}\newtheorem{coro}[theo]{Corollary}\newtheorem{conj}[theo]{Conjecture}
\theoremstyle{definition}\newtheorem{defi}[theo]{Definition}\newtheorem{rema}[theo]{Remark}\newtheorem{exam}[theo]{Example}\newtheorem{exem}[theo]{Example}\newtheorem{nota}[theo]{Notation}
\newcommand{\ie}{i.e.}\newcommand{\eg}{e.g.}\newcommand{\resp}{resp.}\newcommand{\skpt}{\par\smallskip\noindent}
\emergencystretch=3em\pagestyle{plain}
\usepackage{subfig}
\usepackage{xpatch}

\DeclareMathOperator{\E}{\mathbf{E}}



\newcommand{\bbN}{\mathbb{N}}
\newcommand{\Z}{\mathbb{Z}}
%\newcommand{\R}{\mathbb{R}}
%\newcommand{\C}{\mathbb{C}}
\newcommand{\T}{\mathbb{T}}
\newcommand{\ind}[1]{\mathbf{1}_{\left\{#1\right\}}}
%\newcommand{\indset}[1]{\mathbf{1}_{#1}}
\newcommand{\floor}[1]{{\left\lfloor #1 \right\rfloor}}
%\newcommand{\ceil}[1]{{\left\lceil #1 \right\rceil}}

\renewcommand{\P}{\mathbf{P}}


\newcommand{\defeq}{:=}

\renewcommand{\S}{\mathcal{S}}
\newcommand{\clH}{\mathcal{H}}
\newcommand{\clF}{\mathcal{F}}
\newcommand{\clG}{\mathcal{G}}
%%%%%%%%%%%%%%%%%%%%%%%%%%%%%%%%%%%%%%%%%%%%%%%%%%%%%%%%%%%
%~ This fixes spacing issues around \left( and \right) and other similar symbols
\let\originalleft\left 
\let\originalright\right
\renewcommand{\left}{\mathopen{}\mathclose\bgroup\originalleft}
\renewcommand{\right}{\aftergroup\egroup\originalright}
%%%%%%%%%%%%%%%%%%%%%%%%%%%%%%
\let\oldin\in
\renewcommand{\in}{\mathchoice{\oldin}{\oldin}{\,\oldin\,}{\,\oldin\,}}
%%%%%%%%%%%%%%%%%Arrow
\renewcommand*{\to}{\mathchoice{\longrightarrow}{\rightarrow}{\rightarrow}{\rightarrow}}
\let\oldmapsto\mapsto
\renewcommand*{\mapsto}{\mathchoice{\longmapsto}{\oldmapsto}{\oldmapsto}{\oldmapsto}}

%%%%%%%%%%%%%%%%%%%%%%%%%%%%%%%%%%%%%%%%%%%%%%%%%%%%%%%%%%%%%%%%%%%%
\graphicspath{{./figures/}}

\newcommand*{\mk}{\mkern -1mu}
\newcommand*{\Mk}{\mkern -2mu}
\newcommand*{\mK}{\mkern 1mu}
\newcommand*{\MK}{\mkern 2mu}

%\hypersetup{urlcolor=purple, linkcolor=blue, citecolor=red}

\newcommand*{\relabel}{\renewcommand{\labelenumi}{(\theenumi)}}
\newcommand*{\romanenumi}{\renewcommand*{\theenumi}{\roman{enumi}}\relabel}
\newcommand*{\Romanenumi}{\renewcommand*{\theenumi}{\Roman{enumi}}\relabel}
\newcommand*{\alphenumi}{\renewcommand*{\theenumi}{\alph{enumi}}\relabel}
\newcommand*{\Alphenumi}{\renewcommand*{\theenumi}{\Alph{enumi}}\relabel}
\let\oldtilde\tilde
\renewcommand*{\tilde}[1]{\mathchoice{\widetilde{#1}}{\widetilde{#1}}{\oldtilde{#1}}{\oldtilde{#1}}}
\let\oldhat\hat
\renewcommand*{\hat}[1]{\mathchoice{\widehat{#1}}{\widehat{#1}}{\oldhat{#1}}{\oldhat{#1}}}

\newcommand{\llbracket}{\mathopen{[\mkern -2.7mu[}}
\newcommand{\rrbracket}{\mathclose{]\mkern -2.7mu]}}
\newcommand{\llbracketclose}{\mathclose{[\mkern -2.7mu[}}
\newcommand{\rrbracketopen}{\mathopen{]\mkern -2.7mu]}}

\newcommand{\llb}{\llbracket}
\newcommand{\rrb}{\rrbracket}
\newcommand{\llbigbracket}{\mathopen{\bigl[\mkern -4.2mu\bigl[}}
\newcommand{\rrbigbracket}{\mathclose{\bigr]\mkern -4.2mu\bigr]}}
\newcommand{\llbg}{\llbigbracket}
\newcommand{\rrbg}{\rrbigbracket}
%%%%%%%%%%%%%%%%%%%%%%%%%%%%%%%%%%%%%%%%%%%%%%%%%%%%%%%%%%%%%%%%%%%%

\makeatletter
\expandafter\def\csname externalRef@fig:IBM\endcsname{2.1}
\LetLtxMacro\NativeRef\ref
\newcommand{\SourceRef}[1]{\ifcsname externalRef@#1\endcsname\href{https://ahl.centre-mersenne.org/item/AHL_2025__8__635_0/}{\csname externalRef@#1\endcsname\ of the source article}\else\NativeRef{#1}\fi}
\AtBeginDocument{\LetLtxMacro\NativeRef\ref\LetLtxMacro\ref\SourceRef}
\makeatother
\begin{document}
\title{Regular-variation tail-power criterion for the deterministic freezing type of the infinite-bin model}
\author{Bastien Mallein \and Sanjay Ramassamy \and Arvind Singh}
\maketitle
\section{Introduction}\label{sec:introduction}

The infinite-bin model (IBM) is a discrete-time particle system on $\Z$ introduced by Foss and Konstantopoulos~\cite{FK}. In this process, at each integer time, a particle chosen according to its rank reproduces by creating a newborn particle one step to its right. Informally, it is constructed as a balls and bins scheme, by placing a bin at each site of $z\in\Z$ with a certain number of balls, corresponding to the particles at initial time. We assume that all the bins at location $z$ large enough are initially empty, so it is possible to number all the balls in the system starting from the rightmost one. We define the dynamics of the infinite-bin model as follows. Let $\mu$ be a probability distribution on $\bbN
:= \{1,2,\ldots\}$. At each integer time, the $k^{\rm th}$ rightmost ball is selected with probability $\mu(k)$ and a ball is added to the bin immediately on its right. See Figure~\ref{fig:IBM} for an illustration of the evolution of the IBM.

We introduce some notation needed to define the infinite-bin model formally. The state space of the IBM is defined as
\[
\S
:= \bigl\{ X : \Z \to \Z_+\cup \{\infty\} \text{ such that } \sup\{ j \in \Z : X(j) \neq 0 \} \in (-\infty,\infty) \bigr\}.
\]
An element $X\in \S$ represents a \emph{configuration of balls} where, for each $k \in \Z$, the bin at position $k$ contains $X(k)$ balls.

Let us point out that the total number of balls in a configuration $X \in \S$ may be finite or infinite (but cannot be zero). We even allow the number of balls in a given bin $k$ to be infinite i.e. $ X(k) = \infty$. For the dynamics of the infinite bin model, such a bin will represent an impassable barrier (the configuration on the left of this bin will not matter and remain unchanged forever). A configuration $X \in \S$ is said to be
\begin{itemize}
\item \emph{locally finite} if $X(k) < \infty$ for all $k\in \Z$,
\item \emph{bounded} if $\sup_{k\in \Z} X(k) < \infty$.
\end{itemize}

Given $\mu$ a probability distribution on $\bbN$ and $X_0 \in \S$ an initial configuration, an infinite-bin model $X$ with reproduction law $\mu$ starting from $X_0$ is constructed as an $\S$-valued Markov process in the following fashion. Given $(\xi_n, n \in \bbN)$ i.i.d. random variables with law $\mu$, at each integer time $n$, $X_n$ is constructed from $X_{n-1}$ by adding a new ball to the bin immediately to the right of the bin containing the $\xi_{n}^{\rm th}$ rightmost ball. In other words, writing $J_{n-1}$ for the index of the bin containing the $\xi_{n}^{\rm th}$ rightmost ball at time $n-1$, defined by
\[
\sum_{j = J_{n-1}}^\infty X_{n-1}(j) \geq \xi_{n} > \sum_{j=J_{n-1} +1}^\infty X_{n-1}(j),
\]
we set\footnote{If such an index does not exist, which is the case if the configuration $X_{n-1}$ contains fewer than $\xi_{n}$ balls, we just set $X_{n}=X_{n-1}$.} $X_{n}(\cdot) = X_{n-1}(\cdot) + \ind{\cdot\,=\,J_{n-1}+1}$.

In contrast, we focus here on a different question and study instead how the number of balls in a given bin grows as time passes. Obviously, for all $k \in \Z$, the number $X_n(k)$ of balls in the bin at position $k$ at time $n$ is non-decreasing in $n$ (since balls are never removed from bins). Hence we can define the random variables
\[
X_\infty(k) = \lim_{n \to \infty}\uparrow X_n(k) \in \Z_+ \cup \{\infty\} \quad \text{a.s.}
\]
We call the event $\{X_\infty(k) < \infty\}$ the \emph{freezing} of bin $k$, as it means that, on this event, there is some finite (random) time after which no ball is ever added to bin $k$. This paper aims to describe conditions on the distribution $\mu$ and the initial configuration $X_0$ for which freezing occurs.

One of the main ways in which bins can have a growing number of balls as $n \to \infty$ stems from the ``cascade effect'': when a bin becomes large, it has a greater chance to be hit by the $(\xi_n)$ and therefore it increases the rate of growth of the bin to it right (while decreasing its own growth rate). This reinforcement scheme can lead to non-freezing of bins for particular initial configurations. As we will see below, it can even be the case that some bins grow to infinity while others freeze a.s. To study the different freezing scenarii in detail, it is useful to consider first the extremal configuration $\widehat{X}_0$ consisting of a single infinite bin (i.e. a barrier) at $0$ while all the other bins are empty:
\begin{equation}\label{def:widehatX}
\widehat{X}_0(k)
:=
\begin{cases}
\infty & \text{ if } k = 0,\\
0 & \text{ otherwise}.
\end{cases}
\end{equation}
We will denote by $(\widehat{X}_n, n \geq 0)$ the IBM process starting from $\hat{X}_0$.

\begin{defi}\label{def:classification}
For $d \in \bbN \cup \{\infty\}$, we say that the probability distribution $\mu$ is of \emph{type $d$} if
\[
d =\inf\left\{k\geq 1 \,:\, \widehat{X}_\infty(k) < +\infty \right\} \quad \text{a.s.}
\]
We say that $\mu$ is of \emph{finite type} (resp. of \emph{infinite type}) if $d <+\infty$ (resp. $d =+\infty$).
\end{defi}

Finding the exact type of a general distribution $\mu$ seems a delicate question. However, with the additional assumption that the distribution has regular variation, one can provide an explicit integral criterion to characterize the type of $\mu$. Recall that a function $\gamma : \Z_+ \to (0,\infty)$ is said to have regular variation with index $\alpha$ (at infinity) if
\[
\lim_{n\to\infty}\frac{\gamma(\floor{u n})}{\gamma(n)} = u^\alpha\quad\text{for all }u>0.
\]
When $\alpha = 0$, the function is said to be slowly varying. The next result shows, in particular, that there exist distributions for any type $d \in \bbN \cup \{+\infty\}$.
\setcounter{theo}{6}
\begin{theo}\label{thm:typeK}
Assume that $(\mu(n), n \geq 1)$ is regularly varying at infinity with index $-\alpha$ for some $\alpha \geq 1$. Let $\bar{\mu}(j) = \mu(\llb j, +\infty\llbracketclose)$ denote its tail. Then $\mu$ has a well-defined type $d$ given by
\[
d = \inf\left\{ h \in \bbN: \; \sum_{j} \bar{\mu}(j)^h < \infty\right\}.
\]
In particular, $\mu$ is of finite type if $\alpha > 1$ and of infinite type if $\alpha=1$.
\end{theo}
\section{General properties of infinite-bin models}\label{sec:preliminaries}

Before studying properties of the infinite-bin model dynamics, we introduce a few extra notations that will allow us to describe in greater details the local dynamic of the model.

Given an infinite-bin model configuration $X \in \S$, we define its \emph{rightmost barrier} as
\begin{equation}\label{def:deltabarrier}
\delta(X)
:= \sup\{j\in \Z : X(j) = +\infty\} \;\in \Z\cup\{-\infty\}
\end{equation}
with the convention $\delta(X) = \sup \emptyset = -\infty$ when $X$ is locally finite. The \emph{front} of $X$ is the location of its rightmost non-empty bin (which is well-defined by definition of a configuration), i.e.
\[
F(X)
:= \max\{ j \in \Z : X(j) \neq 0\} \in \Z.
\]
For all $k \in \bbN$, we also define $B(X,k)$ as the location of the $k^{\rm th}$ rightmost ball in $X$, i.e.
\[
B(X,k)
:= \sup\left\{ j \in \Z : \sum_{i=j}^\infty X(i) \geq k \right\},
\]
again with $\sup \emptyset = -\infty$. We observe that
\[
F(X) = B(X,1)\quad\text{and}\quad\delta(X) = \lim_{k\to\infty} B(X,k).
\]

\begin{rema}
For all $X \in \S$ and $k \in \bbN$, the quantity $B(X, k)$ is well-defined irrespectively of the relative ordering of the balls inside each bin. In other words, the dynamic of the infinite-bin model does not require to fix an ordering of the balls. However, it will be convenient to specify an ordering for the balls in each bin, in order to introduce a genealogical structure to the particle system.
\end{rema}


To give a formal definition of the infinite-bin model dynamics, we introduce the operator $\Phi_\xi : \S \to \S$ that maps a configuration to the new configuration obtained by adding a single new ball immediately to the bin on the right of the $k^{\rm th}$ rightmost ball:
\begin{equation}\label{eq:operatorPhi}
\Phi_\xi(X)
:=
\begin{cases}
\left(X(j) + \ind{j\,=\,B(X,\xi)+1}, j \in \Z\right) & \text{ if } B(X,\xi) > -\infty,\\
X & \text{ otherwise.}
\end{cases}
\end{equation}
Given a probability distribution $\mu$ and an initial configuration $X_0 \in \S$, the infinite-bin model $(X_n,\, n\geq 0)$ with distribution $\mu$ starting from $X_0 \in \S$ can be constructed through the following recursion equation, for $n\geq 0$
\begin{equation}\label{eqn:ibm}
X_{n+1}
:= \Phi_{\xi_{n+1}} (X_n),
\end{equation}
where $(\xi_n, n \geq 1)$ is an i.i.d. sequence of random variables with law $\mu$.


In the rest of this article, we always picture balls as stacked vertically in their bin, and ranked from bottom to top. When a new ball is added, it is placed on the top of the previous balls in the bin. With this ordering, for each ball in $X$, we can speak about
\begin{itemize}
\item its \emph{rank} which corresponds to its index when all balls are ordered starting from the front, from right to left, and then from bottom to top inside each bin.
\item Its \emph{location} which is the index/position of the bin it belongs to.
\item Its \emph{height} which corresponds to its index inside the bin with the balls ordered from bottom to top.
\end{itemize}
Let us note that, by definition, for each $\xi \in \bbN$ the location and heights of the balls already present in $X$ are unchanged in the new configuration $\Phi_\xi(X)$. However, the ranks of the balls in $\Phi_\xi(X)$ are either unchanged or increased by one. See Figure~\ref{fig:IBM} for an illustration of the procedure.
\setcounter{section}{3}
\section{IBM starting with an infinite barrier at \texorpdfstring{$0$}{0}}\label{sec:typeK}


We work in this section under the assumption that $(\mu(n), n \geq 1)$ is a regularly varying sequence, i.e. that there exists $\alpha \geq 1$ such that for all $u > 0$, we have
\begin{equation}\label{eqn:muRegularlyVarying}
\lim_{n \to \infty} \frac{\mu(\floor{n u})}{\mu(n)} = u^{-\alpha}.
\end{equation}
We call $-\alpha$ the index of the regularly varying sequence $(\mu(n), n \geq 1)$. Under this condition, we recall that the function
\begin{equation}\label{eqn:slowlyVarying}
L(n)
:= n^{\alpha} \mu(n)
\end{equation}
is a slowly varying function, i.e. it satisfies $\lim_{n \to \infty} \frac{L(\floor{n u})}{L(n)} = 1$ for all $u > 0$.

We consider $(\hat{X}_n, n \geq 0)$, the IBM($\mu$) starting from the initial condition $\hat{X}_0$ defined by~\eqref{def:widehatX} as the configuration with an infinite number of balls in the bin of index $0$ all other bins empty. We show here that under these conditions, for all $k \in \bbN$, the asymptotic growth rate of $\hat{X}_n(k)$ as $n \to \infty$ is deterministic. This, with a chaining argument, allows us to determine the type of an infinite-bin model. The main result of the section is the following, which is a refinement of Theorem~\ref{thm:typeK}.
\begin{prop}\label{prop:laProp}
Assume that $(\mu(n), n \geq 1)$ is regularly varying with\linebreak index~$-\alpha$. Recall the notation $\bar{\mu}(n)
:= \mu(\llb n,+\infty\llb )$ for the tail of $\mu$. For all $k \in \bbN$, we have
\[
\P\Big(\hat{X}_\infty(k) < +\infty\Big) =
\begin{cases}
0 & \text{if }\sum_{j = 1}^\infty \bar{\mu}(j)^k = \infty,\\
1 & \text{if }\sum_{j = 1}^\infty \bar{\mu}(j)^k < \infty.
\end{cases}
\]
Moreover,
\begin{enumerate}
\item If $\alpha > 1$, then for all $k\geq 1$ such that $\sum_{j = 1}^\infty \bar{\mu}(j)^k = \infty$, we have
\begin{equation*}
\lim_{n \to \infty} \frac{\hat{X}_n(k)}{\sum_{j=1}^n \bar{\mu}(j)^k} = \frac{\Gamma\left(1 + \frac{1}{\alpha - 1}\right)}{\Gamma\left(k + \frac{1}{\alpha -1} \right)} \quad \text{a.s.}
\end{equation*}
\item If $\alpha = 1$, then for all $k \in \bbN$,
\[
\hat{X}_n(k) \underset{n\,\to\,\infty}{\sim} nL^{(k)}(n) \quad \text{a.s}
\]
with $L^{(k)}$ a deterministic slowly varying function.
\end{enumerate}
\end{prop}

We will make use of the following Borel--Cantelli lemma/law of large numbers for sums of Bernoulli random variables.


\begin{lemm}\label{lem:llnber}
Let $(Z_n, n\geq 1)$ be a sequence of Bernoulli r.v. adapted to some filtration $(\clF_n)$. Suppose that there exists a deterministic sequence $(p_n, n\geq 1)$ of positive numbers such that
\[
\lim_{n\to\infty} \frac{\P(Z_{n} = 1\mid \clF_{n-1})}{p_n} = 1 \quad\text{a.s.}
\]
Then, we have
\begin{equation}\label{bernoA}
\P\left(\sum_{k=n}^\infty Z_k < +\infty\right) =
\begin{cases}
1 & \text{if }\sum_{n=1}^\infty p_n < \infty,\\
0 & \text{if }\sum_{n=1}^\infty p_n = \infty.
\end{cases}
\end{equation}
Furthermore, in the case that $\sum_{n=1}^\infty p_n = \infty$, we have
\begin{equation}\label{bernoB}
\lim_{n\to\infty}\frac{\sum_{k=1}^n Z_k}{\sum_{k=1}^n p_k} = 1\quad \text{a.s.}
\end{equation}
\end{lemm}

This result can be found in Freedman~\cite{Freedman}, with the equivalence~\eqref{bernoA} being proved in the main theorem, while~\eqref{bernoB}, which is referred to as an extension of Lévy's strong law~\cite[Chapter~6]{Levy}, is proved in~\cite[(40)]{Freedman}. We present here a short self-contained proof, based on an explicit construction that might have independent interest.

\begin{proof}
Equivalence~\eqref{bernoA} is also a special case of the Borel--Cantelli theorem of~\cite{Chen78}, using that $Z_n$ is $\clF_n$-adapted and
\[
\sup_{n \geq 1} \frac{Z_n}{1 + Z_1 + \cdots + Z_n} \leq 1 \quad \text{a.s.}
\]
We now assume $\sum_k p_k = +\infty$ and prove~\eqref{bernoB}. By considering a possibly enlarged probability space, we can assume that we are also given a sequence $(U_n, n\geq 1)$ of i.i.d. variables with uniform distribution on $[0,1]$ independent of everything else. We define the filtration $\clG = (\clG _n, n\geq 0)$ by $\clG _0 = \sigma(\emptyset)$ the trivial $\sigma$-algebra and $\clG _n = \sigma(\clF_n, (U_k)_{k\,\leq\,n})$. For $n\geq 1$, we set
\[
\tilde{Z}_{n}
:= \ind{U_{n}\,\leq\,\P(Z_{n}\,=\,1\mid \clG _{n-1})}.
\]
Both processes $Z$ and $\tilde{Z}$ are adapted to the filtration $(\clG _n)$ and we have, almost surely,
\[
\P\left(\tilde{Z}_{n} = 1 \, \middle| \, \clG _{n-1}\right) = \P\left(U_{n} \leq \P\left(Z_{n} = 1\,\middle|\, \clG _{n-1}\right) \, \middle| \, \clG _{n-1}\right) = \P\left(Z_{n} = 1 \, \middle| \, \clG _{n-1}\right)
\]
Therefore, the processes $Z$ and $\tilde{Z}$ have the same law so we just need to prove~\eqref{bernoB} for $\tilde{Z}$. By hypothesis, we have $\P(Z_{n} = 1\mid \clG _{n-1}) = \P(Z_{n} = 1\mid \clF_{n-1}) \sim p_n$ almost surely as $n\to\infty$. Thus, if we fix $\varepsilon >0$, we have that $\tilde{Z}_n = \ind{U_{n} \,\leq\,\P(Z_{n}\,=\,1\mid \clG _{n-1})} \leq \ind{U_n\,\leq\, (1+\varepsilon)p_n}$ for all $n$ large enough, almost surely. Hence, we deduce that, because $\sum_k p_k$ diverges,
\[
\limsup_{n\to\infty} \frac{\sum_{k=1}^n \tilde{Z}_k}{\sum_{k=1}^n p_k} \leq \limsup_{n\to\infty} \frac{\sum_{k=1}^n \ind{U_k\,\leq\,(1+\varepsilon)p_k}}{\sum_{k=1}^n p_k} = 1+\varepsilon.\quad\text{a.s.}
\]
where we used~\cite[Theorem~2.3.8]{Durrett2010} to compute the limit on the r.h.s. above since the variables $\ind{U_{k} \,\leq\,(1+\varepsilon)p_k}$ are independent Bernoulli variables. The proof for the liminf is identical and we conclude that~\eqref{bernoB} holds.
\end{proof}



We now turn to the proof of Proposition~\ref{prop:laProp}, splitting it into two parts. We first treat the case when $\mu$ is regularly varying with index $-\alpha < -1$, before turning to the case $-\alpha = -1$. We recall the statement of the result in that case:
\begin{lemm}\label{lem:alphaPos}%
\!\!
Assume that $(\mu(n), n\!\geq\! 1)$ is regularly varying with index $-\alpha\! <\! -1$. For all $k\geq 1$ such that $\sum_{j = 1}^\infty \bar{\mu}(j)^k = \infty$, we have
\begin{equation}\label{eqn:growthRatelemma}
\lim_{n \to \infty} \frac{\hat{X}_n(k)}{\sum_{j=1}^n \bar{\mu}(j)^k} = \frac{\Gamma\left(1 + \frac{1}{\alpha - 1}\right)}{\Gamma\left(k + \frac{1}{\alpha -1} \right)} \quad \text{a.s.}
\end{equation}
whereas $\hat{X}_\infty(k) < \infty$ a.s. whenever $\sum_{j = 1}^\infty \bar{\mu}(j)^k < \infty$.
\end{lemm}
\begin{proof}
As a preliminary, we observe that by Karamata's theorem (see Feller \cite[Chapter~VIII.9, Theorem~1]{Feller71}), the sequence $(\bar{\mu}(j), j \in \bbN)$ is regularly varying with index $1 - \alpha < 0$, and more precisely
\begin{equation}\label{eqn:theSlowlyVarying}
\bar{\mu}(n) \sim \frac{n^{-(\alpha - 1)}}{\alpha - 1} L(n) \quad \text{as }\; n \to \infty,
\end{equation}
with $L$ the slowly varying function defined in~\eqref{eqn:slowlyVarying}. More generally, if $k$ is an integer such that $k(\alpha-1) < 1$, using again Karamata's theorem, we have
\begin{equation}\label{eqn:theSlowlyVaryingGeneralized}
\sum_{j=1}^n \bar{\mu}(j)^k \sim \frac{L(n)^k n^{1 + k(1-\alpha)}}{(\alpha-1)^k(1 + k(1-\alpha))} \quad\text{as }\; n \to \infty
\end{equation}
while the sum converges for $k(\alpha-1) > 1$. In the corner case $k(\alpha-1) = 1$, we have
\begin{equation}\label{eqn:theSlowlyVaryingGeneralized2}
\sum_{j=1}^n \bar{\mu}(j)^k \sim \tilde{L}(n)\quad\text{as }\; n \to \infty
\end{equation}
for some slowly varying function $\tilde{L}$ and the sum may diverge or converge depen\-ding~$\tilde{L}$.


We now prove Lemma~\ref{lem:alphaPos} by iteratively computing the growth rate of bin $k$ until we reach the first index $k$ satisfying
\[
\sum_{j=1}^\infty \bar{\mu}(j)^k < \infty.
\]
Recall that $\hat{X}$ is constructed from the sequence $(\xi_n, n\geq 1)$ of i.i.d. variables with law $\mu$ via~\eqref{eqn:ibm}. We define the filtration $\clF_n = \sigma(\xi_1,\ldots, \xi_n)$. Now, if $\sum_{j=1}^\infty \bar{\mu}(j) = \infty$, applying Lemma~\ref{lem:llnber} (here in the simple case of independent r.v.) we obtain
\begin{equation}\label{eqn:growthRateFirst}
\frac{\hat{X}_n(1)}{\sum_{j=1}^n \bar{\mu}(j)} = \frac{\sum_{j=1}^n \ind{\xi_j\,\geq\,j}}{\sum_{j=1}^n \bar{\mu}(j)} \underset{n\,\to\,\infty}{\to} 1 \quad \text{a.s.},
\end{equation}
proving that $\hat{X}_\infty(1) = \infty$ a.s. and obtaining the growth rate of $\hat{X}_n(1)$ stated in~\eqref{eqn:growthRatelemma}. Note that this result holds irrespectively of any regularity assumption on $(\mu(n))$. Conversely, if $\sum_{j=1}^\infty \bar{\mu}(j) < \infty$, then Lemma~\ref{lem:llnber} states that $\hat{X}_\infty(1) < \infty$ a.s. Thus, we recovered the fact, mentioned in the introduction, that $\mu$ is of type $1$ if and only if it has a first moment.

Next, let $k \geq 2$ such $(k-1)(\alpha-1) < 1$ and assume by induction that~\eqref{eqn:growthRatelemma} holds for all $h \in \llb 1, k-1\rrb $, i.e.
\begin{equation}\label{eqn:recursionHyp}
\lim_{n \to \infty} \frac{\hat{X}_n(h)}{\sum_{j=1}^n \bar{\mu}(j)^{h}} = \frac{\Gamma\left(1 + \frac{1}{\alpha - 1}\right)}{\Gamma\left(h + \frac{1}{\alpha - 1}\right)} \quad \text{a.s.}
\end{equation}
We deduce from~\eqref{eqn:theSlowlyVaryingGeneralized} that, for all $h \in \llb 1, k-1\rrb$,
\begin{equation}\label{eqn:recursionHyp2}
\hat{X}_n(h) \sim \frac{\Gamma\left(1 + \frac{1}{\alpha - 1}\right)}{\Gamma\left(h + \frac{1}{\alpha - 1}\right)} \frac{L(n)^h\, n^{1 + h(1-\alpha)}}{(\alpha-1)^h\,(1 + h(1-\alpha))} \quad \text{a.s. as $n \to \infty$}.
\end{equation}
In particular, $\hat{X}_n(h)$ is negligible compared to $\hat{X}_n(h-1)$ as $n\to\infty$ and we obtain that
\[
R_n
:= \sum_{h=1}^{k-1} \hat{X}_n(h) \sim \hat{X}_n(1) \quad \text{a.s. as $n \to \infty$}.
\]
Now, we observe that
\[
\P\left(\hat{X}_{n+1}(k) = \hat{X}_n(k) + 1\,\middle|\, \clF_n\right) = \mu\left(\llbg n - R_n+1, n - R_n + \hat{X}_n(k-1)\rrbg\right)
\]
as, at time $n$, there are exactly $n-R_n$ balls to the right of bin $k-1$, thus the balls in that bin are ranked between $n - R_n+1$ and $n - R_n + \hat{X}_n(k-1)$. We have the equivalences
\[
n - R_n + \hat{X}_{n}(k-1) \; \sim \; n - R_n + 1\; \sim \; n \quad \text{a.s. as $n \to \infty$}.
\]
Because $\mu$ is regularly varying, for all $\epsilon > 0$, there exists $\delta > 0$ such that
\[
(1 - \delta) \leq \liminf_{n \to \infty} \inf_{m \in [n(1-\epsilon),\,n(1+\epsilon)]} \frac{\mu(m)}{\mu(n)} \leq \limsup_{n \to \infty} \sup_{m \in [n(1-\epsilon),\,n(1+\epsilon)]} \frac{\mu(m)}{\mu(n)} \leq (1+\delta).
\]
Hence, we deduce that $\mu(\llb n - R_n+1, n - R_n + \hat{X}_n(k-1)\rrb) \sim \mu(n)\hat{X}_n(k-1)$ a.s. which means that
\begin{equation*}
\lim_{n \to \infty} \frac{\P\left(\hat{X}_{n+1}(k) = \hat{X}_n(k) + 1\,\middle| \,\clF_n\right)}{\mu(n)\hat{X}_n(k-1)} = 1 \quad \text{a.s.}
\end{equation*}
We note that, by~\eqref{eqn:slowlyVarying}, \eqref{eqn:theSlowlyVarying} and~\eqref{eqn:recursionHyp2}, we have, almost surely, as $n \to \infty$,
\begin{align*}
\mu(n) \hat{X}_n(k-1) &\sim n^{-\alpha} L(n) \frac{\Gamma\left(1 + \frac{1}{\alpha - 1}\right)}{\Gamma\left(k-1 + \frac{1}{\alpha - 1}\right)} \frac{n^{1 + (k-1)(1-\alpha)}L(n)^{k-1}}{(\alpha-1)^{k-1}\bigl(1 + (k-1)(1-\alpha)\bigr)} \\
&\sim \frac{\Gamma\left(1 + \frac{1}{\alpha - 1}\right)}{\Gamma\left(k + \frac{1}{\alpha - 1}\right)} \bar{\mu}(n)^k.
\end{align*}
As a result,
\begin{equation*}
\lim_{n \to \infty} \frac{\P\left(\hat{X}_{n+1}(k) = \hat{X}_n(k) + 1\,\middle|\, \clF_n\right)}{\bar{\mu}(n)^k} = \frac{\Gamma\left(1 + \frac{1}{\alpha - 1}\right)}{\Gamma\left(k + \frac{1}{\alpha - 1}\right)} \quad \text{a.s.}
\end{equation*}
Invoking Lemma~\ref{lem:llnber} with the sequence $Z_n
:= \hat{X}_{n+1}(k) - \hat{X}_n(k)$, we find that, almost surely
\[
\hat{X}_\infty(k) = \sum_{n=0}^{\infty} \left(\hat{X}_{n+1}(k) - \hat{X}_n(k)\right) =
\begin{cases}
+ \infty & \text{if }\sum_{n=1}^\infty \bar{\mu}(n)^k = \infty,\\
< \infty & \text{if }\sum_{n=1}^\infty \bar{\mu}(n)^k < \infty.
\end{cases}
\]
Furthermore, when $\sum_{n=1}^\infty \bar{\mu}(n)^k = \infty$, we obtain the asymptotic
\[
\lim_{n \to \infty} \frac{\hat{X}_n(k)}{\sum_{j=1}^n \bar{\mu}(j)^k} = \frac{\Gamma\left(1 + \frac{1}{\alpha - 1}\right)}{\Gamma\left(k + \frac{1}{\alpha - 1}\right)} \quad \text{a.s.}
\]
hence~\eqref{eqn:recursionHyp} holds for $h=k$. This completes the proof of the induction step.

It remains only to treat the special case $k(\alpha -1) = 1$. We have proved above that, in this case, either $\sum_{j=1}^\infty \bar{\mu}(j)^k < \infty$ and then $\hat{X}_\infty(k) < \infty$ a.s. or $\sum_{j=1}^\infty \bar{\mu}(j)^k = \infty$ and, in view of~\eqref{eqn:theSlowlyVaryingGeneralized2},
\[
\hat{X}_n(k) \sim \sum_{j=1}^n \bar{\mu}(j)^k \sim \tilde{L}(n) \quad \text{a.s. as }\; n\to\infty.
\]
which diverges to infinity. In that second case, we repeat the induction step one last time to find now that $\P(\hat{X}_{n+1}(k+1) = \hat{X}_n(k+1) + 1\mid\clF_n) \sim \mu(n)\tilde{L}(n)$. Finally, because $\mu(n)$ is regularly varying with index $-\alpha < -1$, we have $\sum_{j=1}^\infty \mu(j)\tilde{L}(j) < \infty$ from which we conclude, using Lemma~\ref{lem:llnber} one last time, that $\hat{X}_\infty(k+1) < \infty$ a.s. and the proof of Lemma~\ref{lem:alphaPos} is now complete.
\end{proof}
Finally, its remains to prove Proposition~\ref{prop:laProp} when $\alpha = -1$. We recall the statement of the result in the lemma below.
\begin{lemm}\label{lem:alphaNul}
Assume that $(\mu(n), n \geq 1)$ is regularly varying of index $-1$. For all $k \in \bbN$, we have $\hat{X}_\infty(k) = \infty$ a.s. and furthermore
\[
\hat{X}_n(k) \underset{n\,\to\,\infty}{\sim} nL^{(k)}(n) \quad \text{a.s}
\]
with $L^{(k)} : n \mapsto (n \mu(n))^{k-1} \bar{\mu}(n)$ a slowly varying function.
\end{lemm}

\begin{proof}
The proof is based on an identical method as in Lemma~\ref{lem:alphaPos}. We show by induction that for all $k \in \bbN$, $\hat{X}_n(k)$ is diverging at an almost sure rate which is regularly varying with index~$1$. Note that as $\mu(n)$ is regularly varying of index $-1$, the function $\bar{\mu}(n)$ is now slowly varying.

We first recall that we proved~\eqref{eqn:growthRateFirst} without any assumption on the regularity of $\mu$, i.e.
\[
\hat{X}_n(1) \sim \sum_{j=1}^n \bar{\mu}(j) \quad \text{a.s. as }n \to \infty,
\]
therefore by Karamata's theorem, we now have $\hat{X}_n(1) \sim n \bar{\mu}(n)$ as expected, with $L^{(1)}(n) = \bar{\mu}(n)$.

Let $k \geq 2$, we assume by induction that $\hat{X}_n(k-1) \sim nL^{(k-1)}(n)$ a.s. as $n \to \infty$. Then with similar computations as in the proof of Lemma~\ref{lem:alphaPos}, we obtain
\[
\P\left(\hat{X}_{n+1}(k) = \hat{X}_n(k) + 1\,\middle|\, \clF_n\right) \sim \hat{X}_n(k-1) \mu(n) \quad \text{a.s.}
\]
Note that $\hat{X}_n(k-1)\mu(n) \sim (n \mu(n)) L^{(k-1)}(n) = L^{(k)}(n)$ is slowly varying as a product of slowly varying function. Therefore, using Lemma~\ref{lem:llnber} and Karamata's theorem, we deduce that $\hat{X}_n(k) \sim n L^{(k)}(n)$ as $n\to \infty$, which proves the induction hypothesis as the next step. Remark that in particular $\hat{X}_n(k)$ diverges to $\infty$ almost surely.
\end{proof}
\begin{thebibliography}{ČS24}
\bibitem[FK03]{FK} Sergey Foss and Takis Konstantopoulos, Extended renovation theory and limit theorems for stochastic ordered graphs, Markov Process. Relat. Fields 9 (2003), no. 3, 413–468.
\bibitem[Fre73]{Freedman} David Freedman, Another note on the Borel–Cantelli lemma and the strong law, with the Poisson approximation as a by-product, Ann. Probab. 1 (1973), 910–925.
\bibitem[Lév37]{Levy} Paul Lévy, Théorie de l’addition des variables aléatoires, Monographies des Probabilités; Calcul des probabilités et ses applications, vol. 1, Gauthier-Villars, 1937.
\bibitem[Che78]{Chen78} Louis H. Y. Chen, A short note on the conditional Borel–Cantelli lemma, Ann. Probab. 6 (1978), 699–700.
\bibitem[Dur10]{Durrett2010} Rick Durrett, Probability: Theory and Examples, 4th ed., Cambridge Series in Statistical and Probabilistic Mathematics, vol. 31, Cambridge University Press, 2010.
\bibitem[Fel71]{Feller71} William Feller, An introduction to probability theory and its applications. Vol. II., 2nd ed., Wiley Series in Probability and Mathematical Statistics, John Wiley \& Sons, 1971.
\end{thebibliography}
\end{document}
